% ===========================================================================
\ReportSection[sec:comparison]{Comparative Analysis}

\ReportSubsection[sec:tco]{Three-year cost}
\TodoAnchor{aoa-tco}

Hardware in this table is at street prices; Dell-configured totals are in
Table~\ref{tbl:dell}. Assumptions (also used in the Business Justification and accepted by
the sponsor for now): operator salary
\$100,000/year, or \$48.08/hour with no overhead multiplier. Upkeep: 2~hr/week
for the ML side (given) plus 1~hr/week for Computer~1 (assumed); a second node
adds 1~hr/week (assumed). Setup: 80--160~hr for Computer~1, 160--320~hr for the
first ML node, and 40--80~hr for a second node (all assumed).

\TableTitle[tbl:tco]{Estimated 3-year cost by configuration (USD)}
{\scriptsize
\begin{xltabular}{\textwidth}{@{}l Y{0.891} Y{1.366} Y{0.891} Y{0.95} Y{0.95} Y{0.952}@{}}
  \toprule
  Config. & Hardware & Setup labor & Annual non-labor opex & Annual upkeep & \textbf{3-yr total} & 3-yr ML only$^\dagger$ \\
  \midrule
  \endhead
  S1 & \$30--40K & 240--480~hr; \$11.5--23.1K & \$4--6.5K & 3~hr/wk; \$7.5K & \textbf{\$76--105K} & \$47--63K \\
  S2 & \$39--50K & 240--480~hr; \$11.5--23.1K & \$5--8K & 3~hr/wk; \$7.5K & \textbf{\$88--120K} & \$59--77K \\
  D1 & \$48--62K & 280--560~hr; \$13.5--26.9K & \$6--10K & 4~hr/wk; \$10K & \textbf{\$110--149K} & \$80--107K \\
  D3 & \$57--72K & 280--560~hr; \$13.5--26.9K & \$7--11.5K & 4~hr/wk; \$10K & \textbf{\$122--163K} & \$92--121K \\
  S4 & \$58--73K & 240--480~hr; \$11.5--23.1K & \$7--11K & 3~hr/wk; \$7.5K & \textbf{\$113--152K} & \$84--109K \\
  D4 & \$66--82K & 280--560~hr; \$13.5--26.9K & \$8--13K & 4~hr/wk; \$10K & \textbf{\$134--178K} & \$104--136K \\
  \bottomrule
\end{xltabular}}
{\footnotesize $^\dagger$ Excludes Computer~1 (about \$29--42K over 3 years, including its hardware, setup, operating cost and 1~hr/wk upkeep), which is needed for storage and licensing under any LLM alternative except the status quo. Non-labor opex covers support, software, power and cooling. 3-year total = hardware + setup labor + 3 $\times$ (non-labor opex + upkeep).}

\ReportSubsection{Scoring}

\TableTitle{Weighted scores (1 = poor, 5 = excellent; illustrative, revise with stakeholders)}
{\small
\begin{xltabular}{\textwidth}{@{}Y{1.0} r r r r r r@{}}
  \toprule
  Criterion (weight) & Alt 0 & Alt 2 & S2 & D1 & D3 & D4 \\
  \midrule
  \endhead
  CUI compliance and control (25\%) & 4 & 3 & 5 & 5 & 5 & 5 \\
  Capability and context (15\%) & 1 & 5 & 4 & 3 & 4 & 4 \\
  Concurrency and scalability (15\%) & 2 & 5 & 2 & 2 & 3 & 4 \\
  Resilience and continuity (15\%) & 3 & 5 & 2 & 4 & 4 & 4 \\
  3-year cost (15\%) & 5 & 3 & 4 & 3 & 3 & 2 \\
  Time to capability (5\%) & 5 & 3 & 3 & 3 & 3 & 3 \\
  Integration and customization (10\%) & 1 & 3 & 5 & 5 & 5 & 5 \\
  \midrule
  \textbf{Weighted score} & \textbf{3.00} & \textbf{3.90} & \textbf{3.70} & \textbf{3.70} & \textbf{4.00} & \textbf{4.00} \\
  \bottomrule
\end{xltabular}}

Interpretation: D3 scores as high as D4 at about \$9K less hardware at street prices, and it keeps a
failover path and a dedicated training node. The compliant cloud scores well on the
LLM alone but is sensitive to usage growth and to CUI contract and risk acceptance,
and it does not replace Computer~1. The scores are judgments, not measurements. D1 appears
in the table but does not meet the stated 12 simultaneous users.

\ReportSubsection{Break-even versus a compliant cloud LLM}
\[
  \text{Break-even (months)} = \frac{\text{On-premises ML upfront cost}}{\text{Monthly cloud cost} - \text{Monthly on-premises ML run cost}}
\]
Computer~1 is excluded because it is needed in either case. For D3 at Dell-configured
prices, the ML upfront cost is about \$64K of hardware (street \$49K $\times$ 1.3) plus
about \$14K of setup labor, or about \$78K, and the ML run cost is about \$1.25K/month
(non-labor opex plus upkeep). If cloud spend reaches \$4K/month, break-even is about
28 months. Averaged over 3 years the on-premises ML side costs about \$3.4K/month, so the
cloud is cheaper if spend stays below that. For 20--30 users a compliant cloud LLM is
therefore plausibly cheaper than D3 on usage cost alone unless usage is heavy. The case
for on-premises rests on keeping CUI inside the boundary, training on the NAS archive
without sending it off-site, no per-token cost growth, and the fact that Computer~1 is
needed anyway. Usage profiling in the pilot should drive this choice.

\ReportSubsection{Capability notes}
\begin{itemize}
  \item Open-weight models in 2026 handle most enterprise tasks (summarization,
    extraction, drafting, scripting) well but may trail the best proprietary models on
    the hardest reasoning and agentic tasks.
  \item At this budget the largest practical models are about 120B parameters
    (nodes with 96--192~GB) and a 4-bit 400B-class mixture of experts on a 4-GPU
    node. Quantization lowers quality slightly; validate against real workflows.
  \item Long-context quality degrades with length; use retrieval and chunking instead of
    relying on raw 128K-token prompts for everything.
\end{itemize}
